On-Policy Distillation (OPD) has emerged as a cornerstone technology in the post-training pipeline of Large Language Models (LLMs). Recent advancements, such as DeepSeek-V4~\cite{deepseekv4},MiMo~\cite{xiao2026mimo} and Qwen-3~\cite{yang2025qwen3}, have demonstrated the potent role of OPD as a successive stage following Supervised Fine-Tuning (SFT) and Reinforcement Learning with Verifiable Reward (RLVR) stages. Thinking Machines Lab~\cite{lu2025onpolicydistillation} reproduces and confirms that providing token-level teacher rewards on the student model’s own trajectories can train strong reasoning models efficiently, especially compared with conventional SFT and RLVR. Inspired by this paradigm, self-distillation frameworks—where a single model serves as both teacher and student—further prove that this approach can effectively generate and leverage learning signals to iteratively enhance model capabilities~\cite{zhao2026self,shenfeld2026self,hubotter2026reinforcement}.

During OPD training, the teacher provides token-level losses on the student-generated trajectory, which serve as dense and informative reward signals for the student model. Importantly, the teacher does not penalize all positions equally. Instead, a small subset of \textit{high-loss tokens} receives much stronger penalties, indicating positions where the student may have entered an undesirable reasoning branch. These tokens therefore become the primary correction points in OPD: by assigning larger optimization pressure to them, the teacher guides the student away from erroneous continuations and toward more reliable reasoning trajectories.

This perspective is consistent with recent findings in RLVR, where successful on-policy training has been shown to rely on a small number of \textit{forking tokens} that determine the direction of reasoning paths~\cite{wangbeyond}. Rather than uniformly changing the entire token distribution, RLVR often produces large distributional shifts only at these decisive branching points, allowing the model to acquire correct reasoning behaviors through sparse but critical updates~\cite{meng2026sparse}. OPD can be viewed as providing a more direct and fine-grained version of this mechanism: the teacher explicitly identifies where the student deviates from the preferred reasoning path, estimates the severity of each deviation through token-level loss, and supplies dense corrective signals that enable faster and more stable training.

Typically, OPD training exhibits an S-shaped progression: following an initial exploration phase, the model enters a period of rapid optimization driven by a substantial number of \textit{high-loss tokens} that provide significant gradient signals~\cite{fu2026revisiting}. As training progresses effectively, these high-loss positions are expected to diminish significantly, leading to a natural deceleration and eventual plateau of the learning curve, at which point very few high-loss tokens should remain. However, our empirical investigations reveal a counterintuitive phenomenon: even after the training has reached a state of apparent stagnation, a stubborn population of ``\textbf{Rock Tokens}''—comprising as much as 6\% of the total—persists with consistently high loss. This unexpected discovery raises a fundamental question regarding the role of these recalcitriant tokens: \textit{Do they still function as critical Pillars that provide essential signals for the model's policy, or are they merely Stumbling Blocks that hinder and complicate the alignment process?}



To understand the nature of Rock Tokens, we pose and answer three core research questions:

\begin{itemize}
    \item \textbf{RQ1: Do Rock Tokens still provide useful learning signals?}
   To answer this, we analyze the gradient dynamics of Rock Tokens compared to learnable high-loss tokens, incorporating statistical features such as token entropy and gradient cosine similarity.  Our analysis suggests that Rock Tokens do not function as ``Learning Pillars''; instead, their gradient contributions are frequently stochastic and antagonistic to the primary optimization direction.

    \item \textbf{RQ2: Are some Rock Tokens indispensable to the student reasoning trajectory?}
    We next examine whether a subset of Rock Tokens persists because they are strongly preferred by the student model and necessary for coherent reasoning. Through targeted ablation studies, we identify Rock Tokens whose removal significantly disrupts reasoning or linguistic coherence. These tokens suggest the existence of \textbf{Semantic Pillars}: persistent high-loss tokens that are not useful optimization targets, but still support the student's generated solution.

    \item \textbf{RQ3: Can harmful Rock Tokens be bypassed to improve OPD efficiency?}
    Finally, we investigate whether the remaining Rock Tokens function as \textbf{Stumbling Blocks} that inject unproductive training signals. We test this by freezing or downweighting selected Rock Tokens during early OPD training and evaluating convergence and final performance. The results motivate a more selective distillation strategy that preserves useful semantic anchors while avoiding token-level supervision that slows or destabilizes learning.
\end{itemize}


